\section{Data transcription and integrity protocol}\label{app:protocol}
Primary numerical sources enter the repository only through a
transcription protocol: tables are transcribed to CSV against the
publisher's text layer, validated against printed derived quantities
where the source provides them (for the kinetic dataset of
\citet{rothbaum1979rohde}, the printed $R \cdot P^{1.5}$ checksum column
validates 33 of 33 retained rows; one flagged row is excluded), and
tagged with an evidence tier: T1 for peer-reviewed primary sources, T2
for conference or institutional material, and T3 for trade press, which
may motivate scenario scale only and is never fitted, a rule enforced
in code. Machine-readable sources
replace manual transcription with automated parser integrity checks
(Section~\ref{sec:validation}).

\section{Reproducibility}\label{app:repro}
A single gate script executes the self-checks of all twelve modules
(94 checks at the time of writing) and refuses a green status unless
every check passes; the repository convention is that no commit is made
against a red gate, and all figures and result tables regenerate from
source and raw transcriptions. All stochastic components are seeded,
including the Sobol sampler, whose default scrambled quasi-random
sequence is pinned for run-to-run reproducibility. The Python
environment is pinned in \texttt{requirements.txt}; principal
dependencies are NumPy/SciPy/pandas, LightGBM~\citep{ke2017lightgbm},
SALib~\citep{herman2017salib}, and PHREEQC via
phreeqpython~\citep{parkhurst2013phreeqc} for the (deferred) speciation
layer.
